\section{50-State Ratio Analysis}
\label{sec:analysis}

\subsection{H:S Ratios Across All 49 Bicameral States}
\label{sec:analysis:ratios}

Table~\ref{tab:ratios-all} shows the $H$, $S$, $\gcd(H,S)$, and
NestSection compatibility for all 49 bicameral states.

\begin{table}[h]
\centering\small
\begin{tabular}{llrrrl}
\toprule
State & Region & $H$ & $S$ & $g=\gcd(H,S)$ & Compatible \\
\midrule
Washington    & W    & 98  & 49  & 49  & Yes (2:1) \\
Minnesota     & MW   & 134 & 67  & 67  & Yes (2:1) \\
Alaska        & W    & 40  & 20  & 20  & Yes (2:1) \\
Nevada        & W    & 42  & 21  & 21  & Yes (2:1) \\
Idaho         & W    & 70  & 35  & 35  & Yes (2:1) \\
Montana       & W    & 100 & 50  & 50  & Yes (2:1) \\
Colorado      & W    & 65  & 35  & 5   & Yes (13:7 spine=5) \\
Oregon        & W    & 60  & 30  & 30  & Yes (2:1) \\
Wyoming       & W    & 60  & 30  & 30  & Yes (2:1) \\
North Dakota  & MW   & 94  & 47  & 47  & Yes (2:1) \\
South Dakota  & MW   & 70  & 35  & 35  & Yes (2:1) \\
Kansas        & MW   & 125 & 40  & 5   & Yes (25:8 spine=5) \\
Iowa          & MW   & 100 & 50  & 50  & Yes (2:1) \\
Missouri      & MW   & 163 & 34  & 1   & \textbf{No} ($\gcd=1$) \\
Wisconsin     & MW   & 99  & 33  & 33  & Yes (3:1) \\
Michigan      & MW   & 110 & 38  & 2   & Yes (55:19 spine=2) \\
Ohio          & MW   & 99  & 33  & 33  & Yes (3:1) \\
Indiana       & MW   & 100 & 50  & 50  & Yes (2:1) \\
Illinois      & MW   & 118 & 59  & 59  & Yes (2:1) \\
Kentucky      & S    & 100 & 38  & 2   & Yes (50:19 spine=2) \\
Tennessee     & S    & 99  & 33  & 33  & Yes (3:1) \\
Virginia      & S    & 100 & 40  & 20  & Yes (5:2) \\
North Carolina & S   & 120 & 50  & 10  & Yes (12:5) \\
South Carolina & S   & 124 & 46  & 2   & Yes (62:23 spine=2) \\
Georgia       & S    & 180 & 56  & 4   & Yes (45:14 spine=4) \\
Florida       & S    & 120 & 40  & 40  & Yes (3:1) \\
Alabama       & S    & 105 & 35  & 35  & Yes (3:1) \\
Mississippi   & S    & 122 & 52  & 2   & Yes (61:26 spine=2) \\
Louisiana     & S    & 105 & 39  & 3   & Yes (35:13 spine=3) \\
Arkansas      & S    & 100 & 35  & 5   & Yes (20:7 spine=5) \\
Oklahoma      & S    & 101 & 48  & 1   & \textbf{No} ($\gcd=1$) \\
Texas         & S    & 150 & 31  & 1   & \textbf{No} ($\gcd=1$) \\
New Mexico    & SW   & 70  & 42  & 14  & Yes (5:3) \\
Arizona       & W    & 60  & 30  & 30  & Yes (2:1) \\
California    & W    & 80  & 40  & 40  & Yes (2:1) \\
Hawaii        & W    & 51  & 25  & 1   & \textbf{No} ($\gcd=1$) \\
Pennsylvania  & NE   & 203 & 50  & 1   & \textbf{No} ($\gcd=1$) \\
New York      & NE   & 150 & 63  & 3   & Yes (50:21 spine=3) \\
New Jersey    & NE   & 80  & 40  & 40  & Yes (2:1) \\
Connecticut   & NE   & 151 & 36  & 1   & \textbf{No} ($\gcd=1$) \\
Massachusetts & NE   & 160 & 40  & 40  & Yes (4:1) \\
Rhode Island  & NE   & 75  & 38  & 1   & \textbf{No} ($\gcd=1$) \\
Vermont       & NE   & 150 & 30  & 30  & Yes (5:1) \\
Maine         & NE   & 151 & 35  & 1   & \textbf{No} ($\gcd=1$) \\
New Hampshire & NE   & 400 & 24  & 8   & Yes (50:3 spine=8) \\
Maryland      & S    & 141 & 47  & 47  & Yes (3:1) \\
Delaware      & A    & 41  & 21  & 1   & \textbf{No} ($\gcd=1$) \\
West Virginia & S    & 100 & 34  & 2   & Yes (50:17 spine=2) \\
\bottomrule
\end{tabular}
\caption{NestSection compatibility for all 49 bicameral states.
  ``Yes'' denotes $g = \gcd(H,S) > 1$; ``No'' denotes $g = 1$.
  41 states are compatible; 8 are not (Missouri, Oklahoma, Texas, Hawaii,
  Pennsylvania, Connecticut, Rhode Island, Maine, Delaware).
  Region codes: W=West, MW=Midwest, S=South, NE=Northeast, A=Atlantic.}
\label{tab:ratios-all}
\end{table}

\subsection{Classification}
\label{sec:analysis:classification}

Of the 49 bicameral states:
\begin{itemize}
  \item \textbf{41 compatible} ($g > 1$): NestSection produces a valid
    nested bicameral map. The 2:1 ratio is the most common (22 states):
    the senate has exactly half as many districts as the house, giving a
    spine of $S$ super-regions each containing 2 house and 1 senate district.
  \item \textbf{8 incompatible} ($g = 1$): House and senate district counts
    are coprime. Missouri ($H=163$, $S=34$, $\gcd=1$), Oklahoma ($101$, $48$),
    Texas ($150$, $31$), Hawaii ($51$, $25$), Pennsylvania ($203$, $50$),
    Connecticut ($151$, $36$), Rhode Island ($75$, $38$), Maine ($151$, $35$),
    Delaware ($41$, $21$).
    These states cannot use NestSection and must use independent
    house and senate maps.
\end{itemize}

The predominance of the 2:1 ratio reflects a common constitutional design:
many state constitutions were drafted with the explicit intention that senate
districts contain exactly two house districts. This design is most prevalent
in western states settled after the Civil War, which drafted constitutions
influenced by the 1870s--1890s reform movements.

The eight incompatible states are geographically and politically heterogeneous:
they include large states (Pennsylvania, Texas), small states (Rhode Island,
Delaware), and politically diverse states. Their coprime ratios are historical
accidents rather than constitutional design: they reflect incremental chamber-size
changes that were not coordinated across chambers.
